% Generated draft: requires mathematical and source review.
\documentclass[11pt]{article}
\usepackage[margin=25mm]{geometry}
\usepackage{fontspec}
\setmainfont{Latin Modern Roman}
\usepackage{amsmath,amssymb,mathtools,unicode-math}
\setmathfont{Latin Modern Math}
\usepackage{enumitem,needspace,xurl,hyperref}
\setlength{\parindent}{0pt}
\setlength{\parskip}{5pt}
\setlength{\emergencystretch}{3em}
\newcommand{\N}{\mathbb{N}}
\newcommand{\Z}{\mathbb{Z}}
\newcommand{\Q}{\mathbb{Q}}
\newcommand{\R}{\mathbb{R}}
\newcommand{\C}{\mathbb{C}}
\newcommand{\F}{\mathbb{F}}
\newcommand{\A}{\mathbb{A}}
\newcommand{\PP}{\mathbb{P}}
\newcommand{\E}{\mathbb{E}}
\newcommand{\eps}{\varepsilon}
\newcommand{\dd}{\,\mathrm{d}}
\begin{document}
\section*{UnsolvedMath ID: 30006502}
\textbf{Source problem code:} catalogue entry
\subsection{Definitions and mathematical statement}
Let \(X\) be a Banach space. In standard terminology, \(X\) is \emph{\(K\)-convex} if its finite Rademacher projections are uniformly bounded: on \(\Omega_n=\{-1,1\}^n\) with uniform probability measure, put \(r_j(\omega)=\omega_j\) and
\[
P_n f=\sum_{j=1}^n r_j\,\E(r_j f).
\]
Then \(K\)-convexity means
\[
\sup_{n\geq1}\|P_n\|_{L^2(\Omega_n;X)\to L^2(\Omega_n;X)}<\infty.
\]
The space \(X\) has the \emph{UMD property} if, for some \(q\in(1,\infty)\), there is \(\beta<\infty\) such that every finite \(X\)-valued martingale difference sequence \((d_j)_{j=1}^n\) in \(L^q\), on any probability space with a filtration, and every choice \(\varepsilon_j\in\{-1,1\}\), satisfy
\[
\left\|\sum_{j=1}^n\varepsilon_j d_j\right\|_{L^q(X)}
\leq \beta\left\|\sum_{j=1}^n d_j\right\|_{L^q(X)}.
\]
For an open set \(\Omega\subseteq\C\), \(H^\infty(\Omega)\) denotes the bounded holomorphic functions on \(\Omega\), equipped with the supremum norm.

\textbf{Unresolved source notation.} The supplied record does not define \(M_{p,d,\theta}^X\), the angle \(\theta_p\), or the region \(\Sigma_\theta\). It identifies \(-L_X\) as an Ornstein--Uhlenbeck operator but does not specify its ambient function space, underlying measure, dimension, normalization, or operator domain. The scalar field of \(X\) and any complexification needed for holomorphic functional calculus are also unspecified. These omissions prevent a fully determined mathematical formulation. The following questions therefore retain the source symbols as unresolved notation.

\textbf{Dimension-free bounds.} For every \(K\)-convex Banach space \(X\), is it true that
\[
\forall p\in(1,\infty)\quad
\forall\theta\in(0,\theta_p)\quad
\sup_{d\in\{1,2,\ldots\}}M_{p,d,\theta}^X<\infty?
\]
The additional endpoint question is whether, for every such \(X\),
\[
\forall p\in(1,\infty)\quad
\sup_{d\in\{1,2,\ldots\}}M_{p,d,\theta_p}^X<\infty.
\]
The requested bounds are uniform in \(d\); no uniformity in \(X\), \(p\), or \(\theta\) is asserted.

\textbf{Functional calculus.} For every UMD Banach space \(X\), is it true that
\[
\forall p\in(1,\infty)\quad
\forall\theta\in(0,\theta_p)\quad
-L_X\text{ has a bounded }H^\infty(\Sigma_\theta)\text{ functional calculus}?
\]
The dependence of the realization of \(-L_X\) on \(p\) must be specified before this question is fully defined. The inequalities \(0<\theta<\theta_p\) are retained exactly as supplied; the sector convention and the meaning of the claimed optimal angle require clarification. No endpoint functional-calculus assertion is explicitly requested. In all questions, \(p=1\), \(p=\infty\), and \(\theta=0\) are excluded.
\subsection{Short English statement}
The source asks whether \(K\)-convexity gives bounds uniform over all positive integer dimensions, for every \(1<p<\infty\) and \(0<\theta<\theta_p\), and possibly also at \(\theta=\theta_p\). It separately asks whether UMD spaces give the stated Ornstein--Uhlenbeck \(H^\infty\) calculus throughout the same strict angle range. Essential notation and operator domains are missing, so these are conditional transcriptions requiring mathematical review.
\subsection{Sources}
\begin{itemize}
\item Supplied local UnsolvedMath registry record, permanent ID \(30006502\), title \emph{Dimension-Free Bounds and Optimal Ornstein-Uhlenbeck Calculus Angles}, exact source problem code \texttt{catalogue entry}. The record attributes the problem to Oberwolfach Reports workshop open-problem discussions; no individual proposer is supplied. The proposed year is recorded as 2025.
\item Supplied source citation: \emph{Recent Developments in SPDEs and BSDEs meet Harmonic and Functional Analysis (2026)}, Oberwolfach Reports, pp.\ 2843--2916, DOI \texttt{10.4171/owr/2025/53}. The literature-source label is \emph{Oberwolfach Report 53/2025}. \url{https://doi.org/10.4171/owr/2025/53}. The raw extracted statement contains a page marker 2897; this has not been independently checked as the precise problem location.
\item Supplied literature source: \emph{Holomorphic functional calculus and vector-valued Littlewood--Paley theory}. \url{https://doi.org/10.4171/JEMS/1430}.
\item Literature assessment dated exactly \texttt{2026-08-22}: the supplied curation labels the problem \texttt{partially\_solved} and reports that a bounded \(H^\infty\) calculus is known for UMD spaces at some suboptimal sector angle, while the requested optimal-angle statement remains open. This is supplied dated curation, not an independently verified assessment. The background expressly requests independent verification of time-sensitive status and cited claims.
\end{itemize}
\end{document}
